\section{Book Allocation Problem}
\label{sec:bs-book-allocation}

\problemstatement{We are given $n$ books, where book $i$ has
$\textit{pages}[i]$ pages, and $m$ students. Each student must be allocated
a \textbf{contiguous} range of books (in the given order), every book must
be allocated to exactly one student, and every student must receive at
least one book. We must choose the allocation to \textbf{minimize the
maximum} number of pages assigned to any single student.}

\begin{patternbox}
\emph{``Minimize the maximum load, items assigned in contiguous order,
divided among a fixed number of workers/students''} is exactly the same
shape as Capacity to Ship Packages within D Days
(Section~\ref{sec:bs-ship-packages}): there, we fixed the number of
``buckets'' indirectly through a day limit $D$ and a per-bucket capacity
candidate; here, the number of buckets is fixed directly as $m$ students.
Once you notice that ``contiguous partition into a fixed number of groups,
minimize the maximum group sum'' is the same template regardless of
whether it is phrased as books, packages, or paint sections (see
Section~\ref{sec:bs-painters-partition}), you can reuse the identical
greedy feasibility check every time.
\end{patternbox}

\subsection*{Understanding the problem carefully}

For a fixed candidate maximum $\textit{maxPages}$, greedily simulate
assigning books to the current student until adding the next book's pages
would exceed $\textit{maxPages}$; at that point, move to the next student.
This yields $\textit{studentsNeeded}(\textit{maxPages})$, the minimum
number of students required if no student may hold more than
$\textit{maxPages}$ total pages. As $\textit{maxPages}$ increases, each
student can absorb more books before being forced to stop, so
$\textit{studentsNeeded}$ is non-increasing in $\textit{maxPages}$.

\begin{ideabox}[Candidate Space and Predicate]
The candidate space is integer page totals in
$[\max(\textit{pages}), \sum(\textit{pages})]$ (a single student must be
able to hold at least the largest individual book, and a single student
holding everything always works). The predicate is
$P(\textit{maxPages}) = (\textit{studentsNeeded}(\textit{maxPages}) \le
m)$, and we want the smallest $\textit{maxPages}$ where $P$ holds -- this
is a ``minimize subject to feasibility'' problem, narrowing the same
direction as Sections~\ref{sec:bs-koko}--\ref{sec:bs-ship-packages}, not
the flipped direction of Aggressive Cows.
\end{ideabox}

\subsection*{Why this is identical in structure to ship packages}

The greedy feasibility check here is justified by precisely the same
exchange argument as Section~\ref{sec:bs-ship-packages}: greedily filling
each student with as many books as the candidate maximum allows, before
moving to the next student, never uses more students than any other valid
contiguous allocation respecting that same maximum, because deferring a
book to a later student only ever makes that later student's job harder,
never easier, and books cannot be reordered.

\subsection*{Algorithm}

\begin{enumerate}
  \item Initialize $\textit{lo} \gets \max(\textit{pages})$,
        $\textit{hi} \gets \sum(\textit{pages})$.
  \item While $\textit{lo} \le \textit{hi}$:
  \begin{enumerate}
    \item $\textit{mid} \gets \textit{lo}+(\textit{hi}-\textit{lo})/2$.
    \item Compute $\textit{studentsNeeded}(\textit{mid})$ via greedy
          simulation.
    \item If $\textit{studentsNeeded}(\textit{mid}) \le m$: feasible; try
          smaller: $\textit{hi} \gets \textit{mid}-1$.
    \item Else: infeasible; try larger: $\textit{lo} \gets \textit{mid}+1$.
  \end{enumerate}
  \item Return \textit{lo}.
\end{enumerate}

\begin{algorithm}[H]
\caption{Book Allocation}
\begin{algorithmic}[1]
\Procedure{BookAllocation}{\textit{pages}, $m$}
  \State $\textit{lo} \gets \max(\textit{pages})$, $\textit{hi} \gets \sum(\textit{pages})$
  \While{$\textit{lo} \le \textit{hi}$}
    \State $\textit{mid} \gets \textit{lo}+(\textit{hi}-\textit{lo})/2$
    \State $\textit{students} \gets \Call{StudentsNeeded}{\textit{pages}, \textit{mid}}$
    \If{$\textit{students} \le m$}
      \State $\textit{hi} \gets \textit{mid} - 1$
    \Else
      \State $\textit{lo} \gets \textit{mid} + 1$
    \EndIf
  \EndWhile
  \State \Return \textit{lo}
\EndProcedure
\end{algorithmic}
\end{algorithm}

\subsection*{Dry run}

\dryrun{Take $\textit{pages} = [12, 34, 67, 90]$, $m = 2$.}

$\max = 90$, $\sum = 203$.

\begin{itemize}
  \item $\textit{lo}=90,\textit{hi}=203$: $\textit{mid}=146$. Greedy:
        student 1 takes $12{+}34{+}67=113$ (next, $90$, would make
        $203>146$, stop); student 2 takes $90$. $\textit{studentsNeeded}=2
        \le 2$. Feasible. $\textit{hi}=145$.
  \item Narrowing continues; eventually the search converges to
        $\textit{lo}=113$: with maximum $113$, student 1 takes
        $12{+}34{+}67=113$ (next, $90$, would make $203>113$, stop);
        student 2 takes $90$. Still $2$ students. Checking $112$: student
        1 can only take $12{+}34=46$ (next, $67$, would make $113>112$,
        stop); student 2 takes $67$ (next, $90$, would make $157>112$,
        stop); student 3 takes $90$ -- that needs $3$ students, exceeding
        $m=2$. So $113$ is indeed minimal.
\end{itemize}

The answer is $113$.

\subsection*{C++ implementation}

\begin{lstlisting}
#include <bits/stdc++.h>
using namespace std;

int studentsNeeded(vector<int>& pages, int maxPages) {
    int students = 1;
    long long currentLoad = 0;

    for (int p : pages) {
        if (currentLoad + p > maxPages) {
            students++;
            currentLoad = 0;
        }
        currentLoad += p;
    }
    return students;
}

int bookAllocation(vector<int>& pages, int m) {
    int lo = *max_element(pages.begin(), pages.end());
    int hi = accumulate(pages.begin(), pages.end(), 0);

    while (lo <= hi) {
        int mid = lo + (hi - lo) / 2;
        if (studentsNeeded(pages, mid) <= m) {
            hi = mid - 1;
        } else {
            lo = mid + 1;
        }
    }
    return lo;
}

int main() {
    vector<int> pages = {12, 34, 67, 90};
    cout << "Minimum maximum pages: "
         << bookAllocation(pages, 2) << endl;
    return 0;
}
\end{lstlisting}

\begin{complexitybox}
\textbf{Time Complexity.}
\[
  O(n \log(\sum(\textit{pages}))),
\]
identical in structure to Section~\ref{sec:bs-ship-packages}.
\textbf{Space Complexity.} $O(1)$ auxiliary space.
\end{complexitybox}

\begin{takeawaybox}
``Partition a sequence into a fixed number of contiguous groups to
minimize the maximum group sum'' is a single template with at least three
names on this sheet: Book Allocation, Capacity to Ship Packages, and
Painter's Partition (Section~\ref{sec:bs-painters-partition}). Learn the
template once -- greedy bucket-filling feasibility check nested inside a
minimize-subject-to-feasibility binary search -- and recognize its
disguises rather than re-deriving it from scratch each time.
\end{takeawaybox}
